\documentclass{article}
\usepackage[utf8]{inputenc}
\usepackage[T1]{fontenc}
\usepackage{geometry}
\geometry{
    a4paper,
    left=2.5cm,
    right=2.5cm,
    top=2.5cm,
    bottom=2.5cm
}
\usepackage{graphicx}
\usepackage{float}
\usepackage{hyperref}
\usepackage[english]{babel}
\usepackage{amsmath}
\usepackage{amssymb}
\usepackage{array}
\usepackage{booktabs}
\usepackage{caption}
\usepackage{subcaption}
\usepackage{enumitem}
\usepackage{fancyhdr}
\usepackage{titlesec}
\usepackage{url}
\usepackage{textcomp}
\usepackage{xcolor}
\usepackage{listings}
\usepackage{multirow}
\usepackage{longtable}
\usepackage{tabularx}
\usepackage{times}

% Set up page style
\pagestyle{fancy}
\fancyhf{}
\rfoot{\thepage}
\renewcommand{\headrulewidth}{0pt}

% Title formatting
\titleformat{\chapter}[display]{\normalfont\bfseries}{\chaptertitle\thechapter}{0pt}{\Huge}
\titlespacing*{\chapter}{0pt}{50pt}{40pt}

% Document info
\title{Template Document}
\author{Academic Support Agent}
\date{\today}

\begin{document}
\maketitle

\renewcommand{\orientadorname}{Orientador:}
\renewcommand{\coorientadorname}{Coorientador:}
\renewcommand{\agradecimentosname}{Agradecimentos}

% ----------------------------------------------------------
% ELEMENTOS PR -TEXTUAIS
% ----------------------------------------------------------
\pretextual

% ---
% Insere Capa, Folha de rosto, Ficha catalogr fica (se inserida)
% e folha de aprova o (se inserida).
% ---
\maketitle


% ---
% Dedicat ria
% ---
\imprimirdedicatoria{Este trabalho   dedicado  s crian as adultas que,\\
   quando pequenas, sonharam em se tornar cientistas.}
% ---

% ---
% Agradecimentos
% ---
\imprimiragradecimentos{
Fa a os agradecimentos  queles que direta ou indiretamente contribu ram para que voc  tivesse obtido  xito. Inclua na sua lista agradecimentos aos  rg os de fomento, quando for o caso.
}
% ---

% ---
% Ep grafe
% ---
\imprimirepigrafe{
		``Sua vida pode ser dividida em dois per odos: antes de agora e a partir de agora.''\\
		(Prof. Obvious Stating)
}
% ---

% ---
% RESUMO e ABSTRACT
% ---

% Resumo em portugu s - as palavras entre chaves s o as palavras-chave do trbalho
\begin{resumo}{Latex. Abntex. Normas USP}

 	 Segundo a \citeonline[3.1-3.2]{NBR6028:2003}, o resumo deve ressaltar o  objetivo, o m todo, os resultados e as conclus es do documento. A ordem e a extens o  destes itens dependem do tipo de resumo (informativo ou indicativo) e do  tratamento que cada item recebe no documento original. O resumo deve ser  precedido da refer ncia do documento, com exce o do resumo inserido no  pr prio documento. (\ldots) As palavras-chave devem figurar logo abaixo do  resumo, antecedidas da express o Palavras-chave:, separadas entre si por  ponto e finalizadas tamb m por ponto.
     
    Para auxili -lo com o latex, o Ap ndice 
  \ref{cap_exemplos} apresenta os resultados dos comandos inclu dos no arquivo ape\_comandos/abntex2-modelo-include-comandos.tex 

\end{resumo}

\selectlanguage{english}

%% p gina com t tulo em Ingl s
\imprimiringlescapa

% Inserir folha de aprova o en Ingl s
\includepdf[pages={1}]{folhaAprovacao.pdf}

% Resumo em ingl s
\begin{abstract}{Latex. Abntex.}
This is the English abstract.
\end{abstract}
% ---

% ---
% inserir lista de ilustra es
% ---
\listailustracoes
% ---

% ---
% inserir lista de tabelas
% ---
\listatabelas
% ---

% ---
% inserir lista de abreviaturas e siglas
% ---
\listasiglas{abrev/Abreviaturas}
% ---

% ---
% inserir o sumario
% ---
\sumario
% ---

% ----------------------------------------------------------
% ELEMENTOS TEXTUAIS
% ----------------------------------------------------------



% ----------------------------------------------------------
% Copyright
% ----------------------------------------------------------
\newpage
\thispagestyle{empty}

\ \

\vspace{10cm}

I hereby certify that I have obtained all legal permissions from the owner(s) of each third-party copyrighted matter included in my thesis, and that their permissions allow availability such as being deposited in public digital libraries.

\vspace{5cm}

\begin{center}
student name and signature
\end{center}


% ----------------------------------------------------------
% Introdu o
% ----------------------------------------------------------
% Content from cap_introducao/introducao would be included here



% ----------------------------------------------------------
% Fundamenta o
% ----------------------------------------------------------
% Content from cap_fundamentacao/fundamentacao would be included here


% ----------------------------------------------------------
% Proposta de pesquisa
% ----------------------------------------------------------
% Content from cap_proposta/proposta would be included here



% ----------------------------------------------------------
% Experimentos e avalia o dos resultados
% ----------------------------------------------------------

% Content from cap_experimentos/experimentos would be included here


% ---
% Finaliza a parte no bookmark do PDF, para que se inicie o bookmark na raiz
% ---
\bookmarksetup{startatroot}% 
% ---

% ---
% Conclus o
% ---
% Content from cap_conclusao/conclusao would be included here



 

% ----------------------------------------------------------
% ELEMENTOS P S-TEXTUAIS
% ----------------------------------------------------------
\postextual

% ----------------------------------------------------------
% Refer ncias bibliogr ficas
% ----------------------------------------------------------


% ----------------------------------------------------------
% Gloss rio
% ----------------------------------------------------------
%
%\glossary

% ----------------------------------------------------------
% Ap ndices
% ----------------------------------------------------------
% ---
% Inicia os ap ndices
% ---
\begin{apendicesenv}
% Imprime uma p gina indicando o in cio dos ap ndices
\partapendices
% ----------------------------------------------------------
% Incluir Ap ndice
% ----------------------------------------------------------
% ----------------------------------------------------------
% Capitulo com exemplos de comandos inseridos de arquivo externo 
% ----------------------------------------------------------
% Content from ape_comandos/abntex2-modelo-include-comandos would be included here


\end{apendicesenv}
% ---

% ----------------------------------------------------------
% Anexos
% ----------------------------------------------------------
% ---
% Inicia os anexos
% ---
\begin{anexosenv}
% Imprime uma p gina indicando o in cio dos anexos
\partanexos
% ---
% Incluir Anexo
% ---
\chapter{Morbi ultrices rutrum lorem.}

%o comando lipsum[] comando serve apenas para incluir texto no documento para efeito de visualiza o do formato.
\lipsum[1-25]
\section{Test}
\lipsum[1-20]


\end{anexosenv}
\end{document}